% ============================================================================
% Stage 9 -- Methods subsection (Nature-style)
% Property-guided reinforcement fine-tuning of ChemBERTa on interaction-
% and randomly masked SMILES
%
% Drop this file into the Methods section of your Nature-format manuscript,
% e.g. with \input{stage9_methods.tex}, and merge references.bib into your
% bibliography resource (\bibliography{...,references}). Written against
% natbib-style numeric \cite commands, which the standard Nature/Springer
% Overleaf templates (sn-jnl, nature.cls with natbib) accept unmodified.
% ============================================================================

\subsection*{Property-guided reinforcement fine-tuning of the generative model}
\label{sec:stage9}

To bias the pretrained masked-language model towards generating valid,
drug-like, synthetically accessible, scaffold-diverging and low-toxicity-risk
molecules, we fine-tuned ChemBERTa~\cite{chithrananda2020chemberta} with a
reinforcement-learning objective computed directly on chemistry-aware,
\emph{post-hoc} scores of its own sampled outputs, rather than on a
token-level cross-entropy target.

\paragraph{Rationale for a policy-gradient objective.}
The properties we wished to optimise -- RDKit-validity, drug-likeness (QED),
synthetic accessibility, scaffold novelty relative to the parent ligand, and
two independent toxicity-risk indicators -- can only be evaluated on the
fully decoded, discrete SMILES string produced after every masked position
has been filled and sampled tokens have been mapped back to characters. This
decoding step is not differentiable, so no gradient can be back-propagated
from a downstream cheminformatics score to the token logits by ordinary
autodiff. We therefore treated masked-token infilling as a stochastic policy
and used the REINFORCE (score-function) gradient
estimator~\cite{williams1992reinforce,suttonbarto2018}: each sampled token's
log-probability is retained with its computation graph intact, the finished
molecule is scored by RDKit outside of autodiff, and a scalar loss is
constructed from the \emph{log-probabilities} so that \texttt{loss.backward()}
legitimately propagates a scoring signal into the trainable weights despite
the intervening non-differentiable chemistry evaluation. The same device is
standard for optimising sequence models against non-differentiable
sequence-level metrics~\cite{ranzato2016mixer,rennie2017scst}.

\paragraph{Parameter-efficient adaptation.}
The ChemBERTa backbone (\texttt{ChemBERTa-zinc-base-v1}, pretrained by
masked-language modelling on 77 million ZINC
SMILES~\cite{chithrananda2020chemberta}) was kept frozen throughout, and
trainable low-rank adaptation (LoRA) matrices~\cite{hu2022lora} were inserted
into the query and value projections of every self-attention block (rank
$r=8$, scaling $\alpha=16$, dropout $0.05$), leaving approximately 1\% of the
model's parameters trainable. Adapters were implemented with the
Hugging Face PEFT library~\cite{mangrulkar2022peft}.

\paragraph{Training data.}
Training pairs $(\tilde{s}, s)$ of a masked SMILES $\tilde{s}$ and its
unmasked parent $s$ were drawn from two upstream sources and, unless noted
otherwise, pooled: (i) molecules subjected to uniform random token-level
masking (Stage 1a), and (ii) co-crystallised ligands from the Protein Data
Bank masked at protein--ligand interaction sites identified by
PLIP~\cite{salentin2015plip} (Stage 1b), where the candidate atom pool for
masking was every interacting heavy atom on interaction-bearing residues
(``attractive'' rows) or every non-interacting atom otherwise
(``non-attractive'' rows). To hold the masking rate constant across sources,
each Stage 1b pool was re-sampled, with a deterministic per-molecule seed, to
$\lfloor p \times n_{\mathrm{tok}} \rfloor$ byte-pair-encoded (BPE) tokens
($p=15\%$ of each molecule's tokens by default; the entire pool was masked
when it contained fewer atoms than this target), and the masked SMILES was
rebuilt from that subsample rather than from the full interaction pool.
Stage 1a pairs were drawn only from the combination file already masked at
the matching rate $p$. For the paired comparison reported below, the Stage 1a
arm was constructed from the \emph{same} parent molecules as Stage 1b, masked
at the same rate but at uniformly random token positions, so that the two arms
differ only in \emph{which} tokens are masked --- interaction-derived versus
random --- and not in the underlying molecule set.

\paragraph{Evaluation-set redundancy control.}
Stage 1b records one row per ligand \emph{instance}, keyed by (PDB entry,
residue name, chain, residue number), so a ligand resolved in many deposited
structures contributes many near-equivalent pairs. This redundancy was severe
in practice: the 334{,}535 instance-level pairs surviving the masking step
corresponded to only 41{,}474 distinct parent molecules (8.1 instances per
molecule), and the twenty most frequently deposited species alone accounted for
45.3\% of all instances --- led by sulfate (14.2\%) and glycerol (11.7\%), that
is, crystallisation additives rather than ligands of interest. Because the
frequently deposited species are also small, instance weighting distorted
masking difficulty as well as chemical composition: mean parent length was 33.5
BPE tokens instance-weighted against 60.3 tokens per distinct molecule, mean
masked positions per attempt 4.44 against 8.15, and 40.0\% of instance-level
attempts required only a single mask fill against 5.8\% of distinct molecules.
We therefore deduplicated the evaluation set to one pair per unique parent,
keyed on the canonical (RDKit) SMILES and retaining the first occurrence in a
deterministic file order so that the pre- and post-fine-tuning comparisons are
drawn over an identical molecule set; all reported property distributions are
computed over distinct molecules. The effect is a reweighting rather than a
change in model behaviour: baseline (non-fine-tuned) RDKit validity is 29.8\%
(99{,}812/334{,}535) counting instances but 18.5\% (7{,}652/41{,}474) counting
molecules, and an independent-fill model in which each of $m$ masked positions
is completed correctly with probability $q$, calibrated to reproduce the
instance-weighted value, predicts 10.0\% for the deduplicated mixture --- so
the shift in masking-difficulty composition more than accounts for the observed
change. The concurrent rise in mean QED (0.468 to 0.529) and synthetic
accessibility score (3.02 to 3.49), and the fall in generations identical to
their parent (15.7\% to 10.4\%), follow from the same reweighting.
Deduplication was applied to the evaluation set only; the reinforcement-learning
training pool remained instance-weighted.

\paragraph{Rollout and sampling.}
Masked positions were filled in a single, parallel decoding step rather than
autoregressively. One forward pass over the fully masked input $\tilde{s}$
produced logits at every \texttt{<mask>} position at once; each position was
then sampled independently from that same pass, with logits divided by a
temperature $T=1.2$ and the draw restricted to, and renormalised over, the
top $k=20$ candidates. Masked positions are therefore \emph{conditionally
independent} given the unmasked context $c(\tilde{s})$ --- no mask observes
what any other mask sampled --- which is the native decoding behaviour of a
masked language model and the same independence assumption made by
non-autoregressive translation~\cite{gu2018nat} and by parallel decoding of
conditional masked language models~\cite{ghazvininejad2019maskpredict}. The
policy log-probability is accordingly the factorised sum
\begin{equation}
\log P_{\theta}(\tilde{s} \to \hat{s}) \;=\; \sum_{i \,\in\, \mathcal{M}} \log P_{\theta}\bigl(\hat{t}_i \mid c(\tilde{s})\bigr),
\label{eq:stage9-logprob}
\end{equation}
over the set $\mathcal{M}$ of masked positions, and the filled sequence was
decoded to a candidate SMILES $\hat{s}$. We note that
Eq.~\eqref{eq:stage9-logprob} is a factorised approximation to the joint
likelihood rather than the joint itself, since a masked language model does
not define an explicit autoregressive factorisation over the filled
positions~\cite{wang2019bertmouth}; the same one-shot procedure is used at
training and at evaluation time, so no comparison in this work is confounded
by a change of decoding strategy.

\paragraph{Composite reward.}
Each rollout was scored against its pre-mask parent $s$ by a six-term
composite in $[0,1]$,
\begin{align}
\mathrm{score}(\hat{s}, s) \;=\;
  &\,w_{\mathrm{valid}}\,\mathbb{1}[\hat{s}\ \mathrm{parses\ in\ RDKit}]
  \;+\; w_{\mathrm{QED}}\,\mathrm{QED}(\hat{s})
  \;+\; w_{\mathrm{SA}}\Bigl(1-\tfrac{\mathrm{SA}(\hat{s})}{10}\Bigr) \notag \\
  &+\, w_{\mathrm{nov}}\bigl(1-T_{\mathrm{c}}(s,\hat{s})\bigr)
  \;+\; w_{\mathrm{alert}}\bigl(1-\mathrm{alert}(\hat{s})\bigr)
  \;+\; w_{\mathrm{tox21}}\bigl(1-\hat{p}_{\mathrm{tox21}}(\hat{s})\bigr),
\label{eq:stage9-score}
\end{align}
where $\mathrm{QED}$ is the quantitative estimate of
drug-likeness~\cite{bickerton2012qed}; $\mathrm{SA}\in[1,10]$ is the
Ertl--Schuffenhauer synthetic accessibility
score~\cite{ertl2009sascore}; $T_{\mathrm{c}}(s,\hat{s})$ is the Tanimoto
coefficient between Morgan (ECFP-like) fingerprints of radius 2 and 2048 bits
computed for the candidate and its pre-mask parent, so that
$(1-T_{\mathrm{c}})$ rewards divergence from -- rather than reconstruction of
-- the original scaffold~\cite{rogers2010ecfp}; $\mathrm{alert}(\hat{s})\in\{0,1\}$
flags a match against RDKit's built-in PAINS~\cite{baell2010pains} and
Brenk~\cite{brenk2008alerts} structural-alert catalogues, a lightweight,
always-available reactivity/promiscuity filter; and
$\hat{p}_{\mathrm{tox21}}(\hat{s})$ is the aggregated (mean, by default)
sigmoid probability of toxicological activity across a subset of the twelve
nuclear-receptor and stress-response assays of the Tox21
panel~\cite{huang2016tox21,tox21challenge}, predicted by a separately
fine-tuned ChemBERTa sequence-classification head. The two toxicity terms
were kept independent rather than merged into a single value so that the
rule-based filter continues to contribute even when no Tox21 classifier
checkpoint is available, in which case $w_{\mathrm{tox21}}(1-\hat p_{\mathrm{tox21}})$
is defined to contribute exactly zero (not the maximal value) to avoid
implicitly rewarding unknown toxicity. Any component that cannot be computed
(e.g.\ an invalid molecule) likewise defaults to its worst-case value of
zero. Default weights were
$(w_{\mathrm{valid}}, w_{\mathrm{QED}}, w_{\mathrm{SA}}, w_{\mathrm{nov}},
w_{\mathrm{alert}}, w_{\mathrm{tox21}}) = (0.20, 0.15, 0.10, 0.15, 0.15,
0.25)$.

\paragraph{Anchoring to the pretrained distribution.}
A score-maximising objective alone places no constraint on how far the policy
may move from the chemistry ChemBERTa learned in pretraining, and unconstrained
property optimisation is a well-documented route to degenerate, reward-hacked
molecules~\cite{renz2019failure,gao2020synthesizability}. We therefore added a
Kullback--Leibler control term against the frozen pretrained
model~\cite{jaques2017sequencetutor,ziegler2019finetuning}. Because only the
LoRA matrices are trainable, the reference policy $\pi_{\mathrm{ref}}$ is
obtained by disabling the adapter at inference time, which recovers the
pretrained weights exactly and requires neither a second copy of the model nor
a separate checkpoint that could drift out of synchrony. At every masked
position the exact KL divergence over the full vocabulary $V$ was computed
analytically,
\begin{equation}
D_{\mathrm{KL}} \;=\; \sum_{i \,\in\, \mathcal{M}} \; \sum_{v \,\in\, V} \pi_{\theta}\bigl(v \mid c(\tilde{s})\bigr) \, \log \frac{\pi_{\theta}\bigl(v \mid c(\tilde{s})\bigr)}{\pi_{\mathrm{ref}}\bigl(v \mid c(\tilde{s})\bigr)},
\label{eq:stage9-kl}
\end{equation}
rather than estimated from the sampled token, so that the anchor contributes
no additional Monte-Carlo variance to an already high-variance gradient and is
differentiable directly, without passing through the score-function estimator.
Eq.~\eqref{eq:stage9-kl} is evaluated on the raw logits, deliberately not on
the temperature-scaled and top-$k$-truncated distribution actually sampled
from: the anchor is intended to constrain the \emph{model}, so that altering
the sampling policy does not silently redefine what counts as drift. Two
implementation details follow from the same requirement. The reference pass is
evaluated deterministically, and dropout in the frozen backbone --- which
regularises no trainable parameter --- was disabled during Stage 9 training,
because a freshly initialised adapter is by construction identical to the
reference (LoRA is initialised with $B=0$) yet, measured across stochastic
forward passes, would otherwise register $0.1$--$0.3$ nats per position of
apparent divergence that no weight update had caused. Dropout within the LoRA
branch itself is retained.

\paragraph{Optimisation.}
The complete loss for a rollout was
\begin{equation}
\mathcal{L} \;=\; -\bigl(\mathrm{score}(\hat{s},s) - b\bigr)\,\log P_{\theta}(\tilde{s}\to\hat{s}) \;+\; \beta \, D_{\mathrm{KL}},
\label{eq:stage9-loss}
\end{equation}
with $\beta = 0.05$ and using an exponential moving-average reward baseline $b$
(decay $0.99$, updated from the mean score of each mini-batch) as a
variance-reduction control variate~\cite{greensmith2004variance,mnih2014nvil},
with $b$ initialised to zero. Setting $\beta=0$ recovers the unanchored
objective. Per-example losses were averaged over mini-batches of 16 pairs, and
adapter weights were updated with Adam~\cite{kingma2015adam} (learning rate
$5\times10^{-5}$) after clipping the global gradient norm to 1.0. Training ran
for 10 epochs over the pooled, shuffled Stage 1a/1b pairs, with the LoRA
adapter, optimiser state and full training history checkpointed after every
epoch to allow interruption and exact resumption. The mean KL per masked
position was recorded each epoch alongside the reward components, providing a
direct read-out of whether $\beta$ is holding the policy near its pretrained
distribution.

\paragraph{Hardware utilisation.}
The objective above is kept deliberately separate from the way it is executed,
so that the identical estimator can be run either as a per-molecule reference
implementation or through a throughput-oriented path that adapts to whichever
device a run lands on. In the latter, each mini-batch is right-padded into a
single $[B, L]$ tensor and costs exactly one policy forward pass and one
reference forward pass; all $M$ masked positions in the batch are drawn with a
single \texttt{multinomial} call over the gathered $[M, k]$ top-$k$ matrix, and
the per-molecule log-probabilities of Eq.~\eqref{eq:stage9-logprob} are
recovered by a scatter-add over the flat mask index, removing the loop over
positions and the loop over molecules together. Padding cannot contaminate real
positions, since the attention mask enters as an additive $-\infty$ bias and
position indices are a cumulative count over non-padding tokens; batched and
unbatched logits agree to $\sim\!6\times10^{-6}$ at real positions, a
float-associativity difference rather than leakage. Because tokenised SMILES
here span roughly 8 to 500 BPE tokens, pairs are length-bucketed before batching
--- optionally packed to a fixed padded-token budget rather than a fixed
molecule count --- so that one long molecule does not pad an entire batch out to
its length, while batch order is reshuffled every epoch to keep the gradient
sequence stochastic~\cite{ott2018scaling}. Matrix multiplications run under
bfloat16 autocast on Ampere-class and later accelerators, and float16 with
gradient scaling on Volta/Turing, with TF32 enabled for the fp32 matmuls that
remain~\cite{micikevicius2018mixed}; the log-softmax, the sampled
log-probabilities and the KL of Eq.~\eqref{eq:stage9-kl} are always evaluated in
fp32, being the numerically delicate terms and negligible in cost beside the
forward pass. On the host side, RDKit scoring --- the bottleneck once the
forward pass is batched, at roughly $9$\,ms per molecule against $\sim\!1$\,ms
of amortised forward --- is dispatched to a persistent process pool whose
workers each pin their BLAS threads to one, while the Tox21 term is excluded
from that pool and evaluated on the accelerator as one batched classifier
forward so that no worker holds a second copy of the network. Worker and thread
counts follow the job's actual allocation (scheduler variables, then the process
CPU-affinity mask, then the cgroup quota) rather than the host core count, which
on a shared node would oversubscribe cores the job was never granted.
Multi-device training uses DistributedDataParallel~\cite{li2020pytorchddp} with
the configured batch treated as a \emph{global} batch split across ranks, so the
number of optimiser updates matches a single-device run and the gain is purely
in wall-clock time; the baseline $b$ is all-reduced so that ranks share one
control variate, ranks receive an equal number of batches so that none leaves
the gradient all-reduce early, and checkpoints are written by rank 0 alone.
Batch size is accordingly set from the optimiser-step budget rather than from
free device memory: at 44 million parameters the backbone cannot fill a modern
accelerator, and enlarging the batch until it did would only reduce the number
of high-variance REINFORCE updates taken over the same data. Batched sampling
consumes the random-number stream in a different order from per-molecule
sampling, so the two execution paths are compared distributionally rather than
by exact output matching, with token-for-token equivalence asserted under the
deterministic $k=1$ setting.

This overall strategy --- pretrain a generative model on a large SMILES
corpus, then steer it towards desired properties with a policy-gradient
objective --- follows established practice in de novo molecular
design~\cite{segler2018focused,olivecrona2017reinvent,popova2018release,guimaraes2017organ,blaschke2020reinvent2},
and REINFORCE specifically has been examined against alternative optimisers in
this setting~\cite{neil2018exploring}. The formulation used here differs from
those precedents in that the policy is a masked, non-autoregressive infilling
model over a fixed scaffold of unmasked context rather than an autoregressive
generator sampling molecules from scratch.

% ------------------------------------------------------------------------
% Optional hyperparameter summary table -- delete if the journal's Methods
% word limit does not accommodate a table, or move it to the Supplementary
% Information.
% ------------------------------------------------------------------------
\begin{table}[htbp]
\centering
\caption{Stage 9 default hyperparameters.}
\label{tab:stage9-hparams}
\begin{tabular}{@{}ll@{}}
\toprule
Parameter & Value \\
\midrule
Base model & ChemBERTa-zinc-base-v1~\cite{chithrananda2020chemberta} \\
LoRA rank $r$ / $\alpha$ / dropout & 8 / 16 / 0.05 \\
LoRA target modules & attention query, value projections \\
Mask rate $p$ (\% of BPE tokens) & 15\% \\
Decoding & one-shot (all masks in a single forward pass) \\
Sampling top-$k$ / temperature $T$ & 20 / 1.2 \\
KL anchor weight $\beta$ & 0.05 \\
Reference policy $\pi_{\mathrm{ref}}$ & same model, LoRA adapter disabled \\
Backbone dropout during Stage 9 & disabled (LoRA dropout retained) \\
Batch size & 16 \\
Optimiser / learning rate & Adam / $5\times10^{-5}$ \\
Gradient-norm clip & 1.0 \\
Baseline decay & 0.99 \\
Epochs & 10 \\
Reward weights $(w_{\mathrm{valid}}, w_{\mathrm{QED}}, w_{\mathrm{SA}}, w_{\mathrm{nov}}, w_{\mathrm{alert}}, w_{\mathrm{tox21}})$ & (0.20, 0.15, 0.10, 0.15, 0.15, 0.25) \\
\bottomrule
\end{tabular}
\end{table}


% ============================================================================
% Stage 10 -- Methods subsection (Nature-style)
% Property-guided SUPERVISED fine-tuning of ChemBERTa by best-of-K selection.
% Companion to the Stage 9 subsection above: same data, same masking rate,
% same rollout distribution and the same property measurement code; only the
% estimator (and the trainable-parameter budget) differ.
% ============================================================================

\subsection*{Property-guided supervised fine-tuning by best-of-$K$ selection}
\label{sec:stage10}

The reinforcement objective of Eq.~\eqref{eq:stage9-loss} is one of two ways to
turn a non-differentiable chemistry score into a weight update. We implemented
the second as a controlled comparison. Where Stage 9 differentiates a
\emph{score-weighted log-probability}, Stage 10 uses the score only to
\emph{choose a target}, and then performs the ordinary token-level
cross-entropy that pretrained the model in the first place. Everything that is
not the estimator is held fixed between the two stages: the training pairs and
their collection order, the masking rate $p=15\%$, the one-shot rollout and its
sampling distribution, and the RDKit measurement code used both for the
objective and for the reported property distributions.

\paragraph{Rationale: selection in place of the score-function estimator.}
The obstruction described above --- validity, QED, synthetic accessibility,
novelty and the structural-alert filters are computed by a cheminformatics
parser acting on a decoded string, so a score returned by RDKit carries no
gradient --- is a property of the \emph{scoring} function, not of the
\emph{training} objective. It can therefore be circumvented without a
score-function estimator at all: sample several completions of the same masked
molecule, let the score decide which one is best, and fit the model to that
completion by maximum likelihood. The chemistry then enters only through a
discrete $\arg\min$, and every quantity that reaches \texttt{backward()} is a
log-probability produced by the network itself. This is expert
iteration~\cite{anthony2017expert} in its rejection-sampling form, also known as
self-taught or reward-ranked fine-tuning~\cite{zelikman2022star,dong2023raft},
and in molecular design it is the hill-climbing fine-tuning used to build
focused libraries from a pretrained SMILES
generator~\cite{segler2018focused,brown2019guacamol}.

Two consequences distinguish it from Eq.~\eqref{eq:stage9-loss}. First, the
direction of the gradient is bounded: the REINFORCE loss contains the term
$-(\mathrm{score}-b)\log P_{\theta}$, which for a below-baseline rollout is a
\emph{positive} multiple of $\log P_{\theta}$ and so drives the model's own
samples towards probability zero, an objective unbounded below and a known
source of degenerate optima; cross-entropy is bounded below by zero and only
ever raises the probability of tokens that were observed to produce a good
molecule. Second, the achievable optimum is capped by the sampler: selection
can only ever imitate the best member of a $K$-sample pool and cannot reach a
completion the model would not already propose, whereas a policy gradient
places no such ceiling on the direction of the update. Stage 10 accordingly
trades reach for stability.

\paragraph{Composite loss.}
Each candidate $\hat{s}$ was scored against its pre-mask parent $s$ by a
composite \emph{loss} --- the mirror image of the Stage 9 reward of
Eq.~\eqref{eq:stage9-score}, so that lower is better ---
\begin{equation}
\ell(\hat{s}, s) \;=\;
\begin{cases}
\;w_{\mathrm{QED}}\bigl(1-\mathrm{QED}(\hat{s})\bigr)
\;+\; w_{\mathrm{SA}}\,\mathrm{clip}_{[0,1]}\!\left(\dfrac{\mathrm{SA}(\hat{s})-1}{9}\right)
\;+\; w_{\mathrm{nov}}\,T_{\mathrm{c}}(s,\hat{s})
\;+\; w_{\mathrm{alert}}\,\mathrm{alert}(\hat{s}),
& \hat{s}\ \text{valid},\\[2.2ex]
\;S \;+\; w_{\mathrm{valid}},
& \hat{s}\ \text{invalid},
\end{cases}
\label{eq:stage10-loss}
\end{equation}
where $S \equiv w_{\mathrm{QED}} + w_{\mathrm{SA}} + w_{\mathrm{nov}} +
w_{\mathrm{alert}}$ and the descriptors are those of
Eq.~\eqref{eq:stage9-score}. Three differences from the Stage 9 reward are
deliberate. The synthetic-accessibility term is normalised over the score's
actual range $[1,10]$ as $(\mathrm{SA}-1)/9$ rather than as $\mathrm{SA}/10$,
so that the easiest-to-synthesise molecule is charged exactly zero. The novelty
term penalises the Tanimoto \emph{similarity} $T_{\mathrm{c}}(s,\hat{s})$
directly, which is what makes reconstruction of the parent ---
$T_{\mathrm{c}}=1$ --- the worst attainable value of that term. And the Tox21
term is omitted from Eq.~\eqref{eq:stage10-loss}, so that the objective is
computable from RDKit alone and depends on no auxiliary learned checkpoint;
this holds for the three execution variants of
Eq.~\eqref{eq:stage10-loss} described below, and is reversed in Stage 10.4,
which reinstates the term as a fifth component (see
\emph{Learned toxicity term} below). Any component that RDKit
parses but cannot evaluate is charged its full weight, the pessimistic reading
being the safe one for a criterion whose only role is to rank candidates.
Weights were $(w_{\mathrm{QED}}, w_{\mathrm{SA}}, w_{\mathrm{nov}},
w_{\mathrm{alert}}) = (0.30, 0.20, 0.30, 0.20)$, summing to $S=1$, and
$w_{\mathrm{valid}}=1$.

Validity enters Eq.~\eqref{eq:stage10-loss} not as one weighted term among six
but as a separating margin. Because the four property terms are individually
bounded by their weights,
\begin{equation}
\sup_{\hat{s}\ \mathrm{valid}} \ell(\hat{s},s) \;=\; S
\;<\; S + w_{\mathrm{valid}} \;=\; \ell(\hat{s},s)\big|_{\hat{s}\ \mathrm{invalid}}
\qquad \text{for any } w_{\mathrm{valid}}>0 ,
\label{eq:stage10-ordering}
\end{equation}
so \emph{every} valid candidate outranks \emph{every} invalid one, irrespective
of its drug-likeness, accessibility, novelty or alert status. Validity is thus
a lexicographic gate on the selection rather than a quantity traded off against
the others, and $w_{\mathrm{valid}}$ sets only the size of the margin, never
the ordering. Encoding the same priority by a large multiplicative penalty
would instead have made the selection criterion numerically incomparable with
the cross-entropy it feeds, which is of order $1$--$10$ nats.

\paragraph{Learned toxicity term (Stage 10.4).}
Stage 10.4 reinstates the Tox21 component of Eq.~\eqref{eq:stage9-score} as a
fifth term of the selection loss,
\begin{equation}
\ell'(\hat{s}, s) \;=\; \ell(\hat{s}, s)
\;+\; w_{\mathrm{tox21}}\bigl(1 - \hat{p}_{\mathrm{clean}}(\hat{s})\bigr),
\qquad
\hat{p}_{\mathrm{clean}} \equiv 1 - \hat{p}_{\mathrm{tox21}},
\label{eq:stage104-loss}
\end{equation}
with the invalid branch likewise charged the full $w_{\mathrm{tox21}}$. The
worst attainable value for a valid candidate accordingly rises from $S$ to
$S' \equiv S + w_{\mathrm{tox21}}$ and an invalid candidate scores
$S' + w_{\mathrm{valid}}$, so the separating margin of
Eq.~\eqref{eq:stage10-ordering} remains exactly $w_{\mathrm{valid}}$ and
validity is still a lexicographic gate rather than a traded-off quantity. We
added the term on top of the existing four weights rather than renormalising
them, so that every other term retains its Stage 10 meaning and the per-term
diagnostics remain comparable across variants; only the totals shift.
$w_{\mathrm{tox21}} = 0.20$, matching $w_{\mathrm{alert}}$, so that the
rule-based and learned toxicity indicators enter on equal footing and their
disagreement on generated molecules is itself observable.

The sign convention inverts relative to Stage 9 and the inversion is
deliberate. In the Stage 9 \emph{reward}, an unavailable classifier
contributes exactly zero so as not to reward unverified toxicity; here, in a
\emph{loss}, a candidate the classifier cannot score is charged the full
weight, which is the same pessimistic reading already applied to every other
term of Eq.~\eqref{eq:stage10-loss}. One consequence is worth stating because
it is silent: were no classifier configured at all, that term would take the
identical constant value for all $K$ candidates of a molecule and could
therefore not alter the arg-min of Eq.~\eqref{eq:stage10-argmin}. The run
would be numerically indistinguishable from Stage 10.2 while reporting a
toxicity curve, so the implementation refuses to begin under that
configuration rather than degrade silently.

The classifier is a twelve-output sigmoid head fine-tuned from the same
ChemBERTa backbone as the generator on the Tox21 training
set~\cite{huang2016tox21,tox21challenge}, so that classifier and generator
share a tokenizer and the classifier sees the token stream the generator
actually emits. Of the $7{,}831 \times 12 = 93{,}972$ label cells, $77{,}946$
($82.9\%$) are observed; a blank denotes an assay not run rather than an
inactive result, so the training objective is a \emph{masked} binary
cross-entropy under which unobserved cells contribute neither loss nor
gradient. Per-assay positive rates span $2.9\%$ (NR-PPAR-$\gamma$) to
$16.2\%$ (SR-ARE); each assay's positive class is therefore reweighted by its
training-split negative-to-positive ratio, capped at $50$ so that no single
rare assay dominates the gradient. Without that reweighting the minimiser
attains a near-optimal masked cross-entropy by predicting inactivity
everywhere, which would again render $\hat{p}_{\mathrm{clean}}$ almost
constant and the term inert. The classifier is partitioned by the same
Bemis--Murcko framework rule applied to the generator, so that the two splits
mean the same thing, and its performance is reported as per-assay ROC-AUC over
the held-out frameworks.

Following the Stage 9.1 arrangement, the classifier is excluded from the RDKit
worker pool and evaluated as a single batched forward on the training device.
Stage 10 scores $K$ candidates per molecule rather than one, so two reductions
are applied, both exact because the classifier is frozen and in evaluation
mode: candidates that fail to parse take the invalid branch of
Eq.~\eqref{eq:stage104-loss} and never read the term, so they are not
submitted; and repeated SMILES --- common, since the $K$ completions are drawn
from a top-$k$ truncated distribution --- are scored once and the value
scattered to their positions.

\paragraph{Candidate pool and target selection.}
For each pair $(\tilde{s}, s)$ a \emph{single} forward pass produced the logits
$z_i \in \mathbb{R}^{|V|}$ at every masked position $i \in \mathcal{M}$. That
one pass serves both roles in the algorithm: it defines the proposal
distribution the candidates are drawn from, and it supplies the logits against
which the cross-entropy is later computed, so no second forward pass is
required. $K=16$ completions were drawn independently and identically from the
same temperature-scaled, top-$k$-truncated distribution used in Stage 9,
\begin{equation}
\hat{t}^{(k)}_i \;\sim\; q_{\theta}\bigl(\cdot \mid c(\tilde{s})\bigr)
\;\propto\; \mathrm{softmax}\bigl(z_i / T\bigr)\Big|_{\mathrm{top}\text{-}k},
\qquad k = 1,\dots,K, \quad i \in \mathcal{M},
\label{eq:stage10-sampling}
\end{equation}
with $T=1.2$ and $k=20$, and each was decoded to a candidate SMILES
$\hat{s}^{(k)}$. Sampling was performed under \texttt{detach}: the draw is a
\emph{search} over possible targets and carries no gradient, which is precisely
what distinguishes this estimator from REINFORCE, where the sampled token's
log-probability is retained with its computation graph intact. The pool was
scored by Eq.~\eqref{eq:stage10-loss} and the winner taken as
\begin{equation}
k^{\star} \;=\; \arg\min_{k \in \{1,\dots,K\}} \ell\bigl(\hat{s}^{(k)}, s\bigr),
\qquad
\mathcal{V} \;=\; \mathbb{1}\bigl[\hat{s}^{(k^{\star})}\ \text{is valid}\bigr].
\label{eq:stage10-argmin}
\end{equation}
By Eq.~\eqref{eq:stage10-ordering}, $\mathcal{V}=1$ if and only if at least one
of the $K$ candidates parses.

\paragraph{Parent fallback and its weighting.}
Measured on these data at $p=15\%$, approximately $70\%$ of molecules yield
\emph{no} valid completion in $K=16$ draws, and $K=8$ gives an
indistinguishable rate --- validity is bimodal and molecule-intrinsic here, not
a matter of drawing enough samples. Those molecules offer nothing worth
imitating, so the target reverts to the parent's own tokens, which are
available for every pair and valid by construction:
\begin{equation}
\mathbf{t}^{\star} \;=\;
\begin{cases}
\bigl(\hat{t}^{(k^{\star})}_i\bigr)_{i \in \mathcal{M}}, & \mathcal{V}=1,\\[1ex]
\bigl(t^{\,s}_i\bigr)_{i \in \mathcal{M}}, & \mathcal{V}=0,
\end{cases}
\qquad
\omega \;=\;
\begin{cases}
1, & \mathcal{V}=1,\\[1ex]
\lambda_{\mathrm{fb}}, & \mathcal{V}=0,
\end{cases}
\label{eq:stage10-target}
\end{equation}
where $t^{\,s}_i$ is the token of the canonical parent at position $i$.
Recovering $t^{\,s}_i$ requires no sequence alignment and involves no
heuristic. The masker tokenises the \emph{canonical} SMILES and replaces whole
BPE tokens by single \texttt{<mask>} symbols, so re-tokenising the masked
string cannot change its length: an unmasked run is the same substring and
re-merges identically, and each \texttt{<mask>} is a hard boundary that
prevents neighbouring tokens merging across it. The masked string therefore has
the same token count as its canonical parent and mask position $i$ holds what
was parent token $i$; verified at $100\%$ on 500 molecules, against only $16\%$
agreement when the raw rather than the canonical parent is used. The length
identity is asserted at every step and a pair that violates it is dropped
rather than trained on a silently misaligned target.

The fallback target is nonetheless a reconstruction target, and reconstructing
the parent attains $T_{\mathrm{c}}=1$, the worst possible value of the novelty
term of Eq.~\eqref{eq:stage10-loss}. Left at full weight it would make roughly
seven in ten updates train a copier. The scalar $\lambda_{\mathrm{fb}}$
(default $0.3$) down-weights those updates alone and spans the whole design
space between the two defensible extremes: $\lambda_{\mathrm{fb}}=1$ trains
them as ordinary examples, maximising the rate at which validity is acquired at
the greatest risk of collapse onto the parent, and $\lambda_{\mathrm{fb}}=0$ is
exactly equivalent to discarding those molecules.

\paragraph{The differentiable objective.}
With the target fixed, the gradient-carrying part of Stage 10 is ordinary
masked-language-model cross-entropy on the logits of the original forward pass,
averaged over a molecule's masked positions so that a heavily masked molecule
does not dominate a mini-batch merely by having more of them:
\begin{equation}
\mathcal{L}_{\mathrm{CE}} \;=\; -\frac{1}{|\mathcal{M}|}
\sum_{i \,\in\, \mathcal{M}} \log P_{\theta}\bigl(t^{\star}_i \mid c(\tilde{s})\bigr).
\label{eq:stage10-ce}
\end{equation}
In variant~10a this is supplemented by an unlikelihood
term~\cite{welleck2020unlikelihood} that explicitly suppresses the tokens
observed to produce invalid molecules. Writing
\begin{align}
\mathcal{B} \;&=\; \bigl\{ (k,i) \;:\; \hat{s}^{(k)}\ \text{invalid},\;\;
i \in \mathcal{M},\;\; \hat{t}^{(k)}_i \neq t^{\star}_i \bigr\}, \notag \\
\mathcal{L}_{\mathrm{UL}} \;&=\; -\frac{1}{|\mathcal{B}|}
\sum_{(k,i) \,\in\, \mathcal{B}}
\log\Bigl(1 - P_{\theta}\bigl(\hat{t}^{(k)}_i \mid c(\tilde{s})\bigr)\Bigr),
\label{eq:stage10-ul}
\end{align}
the per-molecule objective is
\begin{equation}
\mathcal{L} \;=\; \omega\,\mathcal{L}_{\mathrm{CE}}
\;+\; \lambda_{\mathrm{UL}}\,\mathcal{L}_{\mathrm{UL}},
\label{eq:stage10-total}
\end{equation}
averaged over the molecules of a mini-batch. Variant~10b is
Eq.~\eqref{eq:stage10-total} with $\lambda_{\mathrm{UL}}=0$, that is,
cross-entropy towards the best candidate alone; the two variants were trained
from identical code and data into separate output directories so that the
contribution of the unlikelihood term is isolated. Three details of
Eq.~\eqref{eq:stage10-ul} are deliberate. The penalty is written on
$1-P_{\theta}$ rather than as $+\log P_{\theta}$ because the latter diverges as
the probability is driven to zero, reintroducing exactly the unbounded
direction that motivated abandoning the policy gradient. Positions at which a
rejected candidate happens to agree with the target are excluded from
$\mathcal{B}$, since including them would set the two terms of
Eq.~\eqref{eq:stage10-total} against each other on the same token. And
$\mathcal{L}_{\mathrm{UL}}$ is \emph{not} scaled by $\omega$: on a molecule for
which no candidate was valid, sixteen decoded strings that failed to parse are
legitimate evidence about what should not be generated, unlike the parent
reconstruction that $\omega$ exists to discount. Gradients flow through
$P_{\theta}$ alone; $\mathbf{t}^{\star}$, $\mathcal{B}$ and every value of
$\ell$ are constants of the step.

\paragraph{Interpretation.}
Informally, the chemistry acts as a teacher rather than as a signal. At each
step the model proposes sixteen ways of completing the same masked molecule;
RDKit reads all sixteen, discards those that are not molecules at all, ranks
what remains by drug-likeness, ease of synthesis, divergence from the parent
and absence of structural alerts, and hands back a single answer. The model is
then taught that answer in exactly the way it was pretrained --- ``at this
masked position, the token was this one'' --- and nothing about the chemistry
needs to be differentiable for that lesson to be delivered, because the only
thing the chemistry ever did was point. Validity improves not because a
validity gradient exists, but because the tokens that are repeatedly pointed at
are the ones whose ring closures balance, whose valences are legal and whose
brackets match. The unlikelihood term of Eq.~\eqref{eq:stage10-ul} is the one
place where the objective speaks about invalidity directly, and it is also the
only signal that survives when the teacher has nothing good to point at.

The same reading explains why Stage 10 carries no analogue of the KL anchor of
Eq.~\eqref{eq:stage9-kl}. The targets of Eq.~\eqref{eq:stage10-target} are
either drawn from the model's own top-$k$ support or are the tokens of a real
deposited molecule, so the objective cannot move probability mass onto
completions the model would never have proposed; the anchoring that
$\beta D_{\mathrm{KL}}$ supplies explicitly in Stage 9 is implicit in the
construction of the target here. The corresponding weakness is the mirror
image: an improvement the sampler never proposes is unreachable, and with a
$70\%$ fallback rate the majority of updates are parent reconstructions at
weight $\lambda_{\mathrm{fb}}$ rather than best-of-$K$ imitations. The
trade-off that $\lambda_{\mathrm{fb}}$ governs --- validity acquired by
reconstruction against novelty lost to copying --- is therefore the central
hyperparameter of this stage, and is read off directly from the diagnostics
below.

\paragraph{Optimisation and trainable parameters.}
Stage 10 fine-tunes ordinary weights rather than adapters: the LM head together
with the last $N=1$ of the six encoder blocks of ChemBERTa-zinc-base-v1 was
left trainable ($7.7$M of $44.1$M parameters, $17.5\%$) and the remainder
frozen. This is a different capacity budget from the Stage 9 adapters, which
touch the query and value projections of \emph{every} block at rank 8, and we
record it here rather than leave it implicit: the two stages differ in their
trainable parameterisation as well as in their estimator, so the comparison
between them is not a clean ablation of the loss alone. Mini-batches of 16
pairs were optimised with Adam~\cite{kingma2015adam} at learning rate
$5\times10^{-5}$ after clipping the global gradient norm to 1.0, for 4 epochs
over the same pooled, shuffled Stage 1a/1b pairs used by Stage 9. Only the
unfrozen tensors are written ($\approx\!30$\,MB against $180$\,MB for the whole
model), together with the optimiser moments and the state of every random
stream, so that an interrupted run resumes as the run it was rather than as a
new one begun from the same weights (see \emph{Checkpointing and run
provenance} below). The trained model was then evaluated with the identical
measurement and figure code used for Stage 9, on the same deduplicated
one-pair-per-parent evaluation set, so that the two stages' property
distributions are directly comparable.

\paragraph{Execution paths and the equivalence between them.}
The objective above is specified once and executed by three interchangeable
implementations, so that the same weights can be produced either by the
simplest arrangement of the computation or by a throughput-oriented one that
adapts to whichever device a run lands on. The reference arrangement performs
one forward pass per molecule and one cheminformatics measurement per
candidate. The second dispatches the $K\!\times\!B$ candidate measurements of a
mini-batch to a persistent process pool whose workers each pin their BLAS
threads to one; because the measurement is a pure function of two decoded
strings, evaluating it in another process cannot alter a floating-point result,
and the pooled and serial objectives agree exactly. This is the speedup that
matters for the present stage: with $K=16$ candidates measured per molecule
against Stage 9's single sampled molecule, the cheminformatics term is roughly
sixteen times more dominant here than there, amounting to $256$ measurements
per optimiser update at the batch size used.

The third arrangement adds the remaining device-level optimisations described
for Stage 9 --- one right-padded $[B,L]$ forward pass per mini-batch, a single
\texttt{multinomial} draw over the gathered top-$k$ matrix for every masked
position in the batch, bfloat16 or float16 autocast with the log-softmax and
the cross-entropy retained in float32, TF32 for the remaining float32 matrix
multiplications, length-bucketed batching, thread and worker counts taken from
the job's actual allocation rather than the host core count, and
DistributedDataParallel~\cite{li2020pytorchddp} with the configured batch
treated as a \emph{global} batch split across ranks so that the number of
optimiser updates is independent of the number of devices.

Unlike the Stage 9 case, these last optimisations are \emph{not} objective
preserving, and the reason is intrinsic to selection rather than incidental to
the implementation. A score that multiplies a log-probability, as in
Eq.~\eqref{eq:stage9-loss}, tolerates a differently drawn sample: the estimator
remains unbiased. A score that \emph{selects} the target of
Eq.~\eqref{eq:stage10-argmin} does not. Three separate effects perturb the
candidate set --- drawing the batch in one \texttt{multinomial} call consumes
the random stream in a different order from one call per molecule; reduced
precision displaces the logits in their trailing bits; and dropout, which is
sampled once per forward pass, supplies $B$ independent masks to $B$ separate
passes but only one to a single padded pass. The third effect is the largest,
displacing the logits at masked positions by $\approx\!1.2$ on this model,
which exceeds the perturbation from reduced precision and cannot be removed by
seeding, since the two arrangements do not request the same shapes from the
generator. Any of the three can change which candidate attains the minimum in
Eq.~\eqref{eq:stage10-argmin}, and the model then learns a different token
sequence rather than a numerically perturbed update. The tuned arrangement is
therefore compared with the reference distributionally --- over mean candidate
validity, mean composite loss and mean novelty --- and never by exact output
matching, while exact agreement is asserted in the deterministic limit of
$k=1$ with dropout disabled, which pins the batched indexing itself.

\paragraph{Checkpointing and run provenance.}
An epoch is not a cheap unit at $K=16$: on the capped Stage 1a/1b pool it is
several hours, so a checkpoint written only at epoch boundaries makes an
interruption expensive and a resumption inexact. The two resumable
implementations therefore write the unfrozen tensors, the Adam moments, the
running per-epoch aggregates, the batch cursor and the state of the Python,
CPU and CUDA random streams every fixed number of optimiser updates, by an
atomic replace so that an interruption during the write cannot destroy the
previous state. Recording the optimiser is not a convenience: rebuilding Adam
after each restart discards the accumulated first and second moments and
imposes a fresh adaptation transient on every resumption, which would otherwise
appear in the training curve as an artefact of the interruption schedule rather
than of the objective. Recording the generator state is what makes a resumed
run a continuation of the run it interrupted, since the same stream draws the
$K$ candidates of Eq.~\eqref{eq:stage10-sampling}.

Because the three implementations optimise one objective, a single run may
optionally be carried across them, beginning under the reference arrangement
and continuing under the tuned one as devices become available. Where that is
done, the checkpoint additionally carries a provenance record naming which
implementation trained which epochs, under which arithmetic precision and
sampling arrangement, and a resumption that changes any of these is reported at
the point it occurs. This is retained deliberately: by the argument above such
a run is not homogeneous in the way either implementation is on its own, the
saved weights carry no trace of the precision under which they were obtained,
and the distinction is not otherwise recoverable after the fact.

\paragraph{Scaffold-based train--validation partition.}
The selection objective of Eq.~\eqref{eq:stage10-loss} is optimised by
cross-entropy towards a target sequence, and on the majority of molecules that
target is the parent's own tokens (the fallback of
Sec.~\ref{sec:stage10}, taken on $\sim$70\% of molecules at $p=15\%$).
Reconstruction is therefore the modal training signal, which makes
memorisation of the training parents the specific failure mode this design
invites, and a property distribution computed over those same parents cannot
distinguish it from generalisation. We therefore partitioned the pooled
training data before any fine-tuning and report all held-out quantities on
molecules the model has never been optimised towards.

The partition is made on \emph{molecular frameworks}, not on instances or on
molecules. Instance-level partitioning is untenable for the reason given above:
a ligand deposited in many structures would appear on both sides of the split
tens of thousands of times. Molecule-level partitioning removes that but leaves
close analogues --- an identical ring system bearing a different substituent
--- free to straddle the split, and such series are abundant in
PDB-derived sets. Partitioning on the Bemis--Murcko
framework~\cite{bemis1996frameworks}, the convention adopted by
MoleculeNet~\cite{wu2018moleculenet}, removes both: no validation molecule
shares a ring system with any training molecule. The consequence is that
held-out figures reported here are systematically lower than the
random-partition figures conventionally quoted; that difference is the
component of apparent performance attributable to memorised chemistry, and its
suppression is the purpose of the exercise. Random partitions are known to
overstate prospective performance for exactly this
reason~\cite{sheridan2013timesplit}, and ligand-based benchmarks constructed
without such control have been shown to reward memorisation over
generalisation~\cite{wallach2018memorization}.

Parents were canonicalised (RDKit~\cite{landrum2023rdkit}) before grouping, so
that a molecule deposited under several equivalent notations forms one group
rather than several that could be assigned to opposite sides. Groups were
ordered by descending size with ties broken on the framework string, and
assigned in that order to validation until the next group would exceed a 10\%
budget, the remainder going to training; every pair inherits the fold of its
parent. Of the 46{,}982 distinct parents recorded in the interaction-masked
source --- counted before the masking-survival filter, and so a superset of
the 41{,}474 quoted above --- 12{,}551 distinct frameworks were obtained. The three largest --- the acyclic
class (6{,}754 molecules, 14.4\%), the Fe-porphyrin/haem framework (4{,}089,
8.7\%) and the Mg-porphyrin/chlorophyll framework (2{,}393, 5.1\%) --- each
exceed the validation budget and are consequently assigned to training, so
validation is drawn from the diverse tail of small frameworks rather than from
the crystallographically over-represented cores. This is the intended
behaviour of a framework partition and not an artefact of the budget. A
further 1{,}015 parents (2.2\%), whose SMILES exceed 400 characters and are
predominantly peptide and oligosaccharide ligands, bypass the framework
reduction and form singleton groups keyed on canonical SMILES; these are
partitioned at molecule granularity, the conservative fallback.

The partition is computed once over the complete parent set and cached as a
manifest read by every Stage 10 variant. It is therefore independent of the
per-parent instance cap, of any sub-sampling of the training pool and of the
variant being run, so a molecule cannot change fold between runs and all
implementations of Sec.~\ref{sec:stage10} agree on what is held out. A
molecule absent from the manifest inherits the fold of its framework, and one
whose framework is also absent defaults to training --- a direction that can
never inflate a held-out score. Validation molecules were evaluated at the end
of each epoch under the identical best-of-$K$ procedure with gradients
disabled; the epoch attaining the lowest mean validation loss was retained as
the final model; and property distributions are reported separately for
training and held-out molecules, the difference between them being the
memorisation measurement that the within-training diagnostics below cannot
supply.

\paragraph{Diagnostics.}
Four quantities were logged per epoch and are the intended read-out of whether
the objective is behaving as constructed: the mean of
Eq.~\eqref{eq:stage10-total}; the mean of $\ell$ at the selected target; the
\emph{candidate validity rate}, the fraction of all $K$ sampled completions
that parse; and the \emph{fallback rate}, the fraction of molecules with
$\mathcal{V}=0$. The last two are the diagnostic pair of interest. A run in
which candidate validity rises while the novelty term of
Eq.~\eqref{eq:stage10-loss} rises alongside it is one that is purchasing
validity by memorising parents, and calls for a smaller
$\lambda_{\mathrm{fb}}$; the intended behaviour is a falling fallback rate with
the novelty term flat or falling.

\begin{table}[htbp]
\centering
\caption{Stage 10 default hyperparameters. Entries shared with Stage 9
(Table~\ref{tab:stage9-hparams}) are held identical by construction.}
\label{tab:stage10-hparams}
\begin{tabular}{@{}ll@{}}
\toprule
Parameter & Value \\
\midrule
Base model & ChemBERTa-zinc-base-v1~\cite{chithrananda2020chemberta} \\
Trainable parameters & LM head $+$ last $N=1$ of 6 encoder blocks (7.7M / 44.1M, 17.5\%) \\
Mask rate $p$ (\% of BPE tokens) & 15\% (shared with Stage 9) \\
Decoding & one-shot (all masks in a single forward pass) \\
Candidates per molecule $K$ & 16 \\
Sampling top-$k$ / temperature $T$ & 20 / 1.2 (shared with Stage 9) \\
Loss weights $(w_{\mathrm{QED}}, w_{\mathrm{SA}}, w_{\mathrm{nov}}, w_{\mathrm{alert}})$ & (0.30, 0.20, 0.30, 0.20), summing to $S=1$ \\
Invalidity margin $w_{\mathrm{valid}}$ & 1.0 (an invalid candidate scores $S+w_{\mathrm{valid}}=2$) \\
Parent-fallback weight $\lambda_{\mathrm{fb}}$ & 0.3 \\
Unlikelihood weight $\lambda_{\mathrm{UL}}$ & 0.5 (variant 10a); 0 (variant 10b) \\
Tox21 term $w_{\mathrm{tox21}}$ & not used (Stages 10, 10.1, 10.2); $0.20$ (Stage 10.4) \\
Worst valid / invalid loss & $S=1$, $S+w_{\mathrm{valid}}=2$; Stage 10.4: $S'=1.2$, $2.2$ \\
Train--validation partition & Bemis--Murcko framework, 10\% held out, seed 42 \\
Model selection & epoch of lowest mean validation loss \\
KL anchor & none (see text) \\
Batch size & 16 \\
Optimiser / learning rate & Adam / $5\times10^{-5}$ \\
Gradient-norm clip & 1.0 \\
Epochs & 4 \\
\bottomrule
\end{tabular}
\end{table}
